% Luc Destombe --- licence CC BY-NC-SA 4.0
% https://creativecommons.org/licenses/by-nc-sa/4.0/deed.fr
% Fichier autonome : compiler avec pdflatex, deux fois (signets, nombre de pages).
\def\lyceeversion{1.0}
\documentclass[a4paper,10pt]{article}
\makeatletter
\newif\iflycee@niveaux
\lycee@niveauxtrue
\newif\iflycee@cm
\newif\iflycee@palatino

\RequirePackage[utf8]{inputenc}
\RequirePackage[T1]{fontenc}
\RequirePackage[french]{babel}
\RequirePackage{etoolbox}

\newcommand{\ShorthandsOff}{\@ifundefined{shorthandoff}{}{\shorthandoff{;:!?}}}
\newcommand{\ShorthandsOn}{\@ifundefined{shorthandon}{}{\shorthandon{;:!?}}}
\ShorthandsOff
\AtBeginDocument{\ShorthandsOn}
\AtBeginEnvironment{tikzpicture}{\ShorthandsOff}

\RequirePackage{geometry}
\RequirePackage{fancyhdr}
\RequirePackage{lastpage}
\RequirePackage{multicol}
\RequirePackage{array}
\RequirePackage{enumitem}
\RequirePackage{xcolor}
\RequirePackage{needspace}

\RequirePackage{amsmath,amssymb}
\iflycee@cm\else\RequirePackage{fourier}\fi
\DeclareMathAlphabet{\mathcal}{OMS}{cmsy}{m}{n}
\renewcommand{\ttdefault}{lmtt}
\begingroup\catcode`\_=\active \catcode`\^=\active
\global\let_\sb \global\let^\sp
\endgroup
\AtBeginDocument{\catcode`\_=12 \mathcode`\_="8000
  \catcode`\^=12 \mathcode`\^="8000 }
\newlength{\retraitcalculs}
\setlength{\retraitcalculs}{2em}

\RequirePackage{tikz}
\usetikzlibrary{arrows.meta,calc,babel,shapes.misc}
\RequirePackage[most]{tcolorbox}
\tcbuselibrary{skins,breakable}

\definecolor{coulDef}{HTML}{1F4E79}
\definecolor{coulProp}{HTML}{7030A0}
\definecolor{coulRem}{HTML}{C55A11}
\definecolor{coulEx}{HTML}{2E7D32}
\definecolor{coulExo}{HTML}{B03A2E}
\definecolor{coulPart}{HTML}{1F4E79}
\definecolor{coulMeth}{HTML}{1B6E4F}
\definecolor{coulCourbe}{HTML}{0B5ED7}
\definecolor{coulCourbeB}{HTML}{C00000}
\definecolor{coulCourbeC}{HTML}{2E7D32}
\definecolor{coulGrille}{HTML}{8C8C8C}
\definecolor{coulRep}{HTML}{1B6E4F}
\definecolor{coulBar}{HTML}{2E7D32}

\newcommand{\R}{\mathbb{R}}
\newcommand{\Rs}{\mathbb{R}^{*}}
\newcommand{\N}{\mathbb{N}}
\newcommand{\Ns}{\mathbb{N}^{*}}
\newcommand{\Z}{\mathbb{Z}}
\newcommand{\Q}{\mathbb{Q}}
\renewcommand{\le}{\leqslant}
\renewcommand{\ge}{\geqslant}
\renewcommand{\approx}{\simeq}
\newcommand{\vect}[1]{\overrightarrow{#1}}
\newcommand{\norme}[1]{\left\lVert #1 \right\rVert}
\newcommand{\intervalle}[2]{\left[#1\,;#2\right]}
\RequirePackage{cancel}
\renewcommand{\CancelColor}{\color{coulExo}}
\newcommand{\fois}[1]{\times{\color{coulExo}#1}}
\newcommand{\barre}[1]{\cancel{#1}}
\newlength{\lycee@eqh}\newlength{\lycee@eqd}
\newsavebox{\lycee@eqa}\newsavebox{\lycee@eqb}\newsavebox{\lycee@eqc}
\newcommand{\lycee@eqcol}[1]{\hbox{$\renewcommand{\arraystretch}{1.35}%
  \begin{array}[t]{@{}>{\displaystyle}r@{}>{\displaystyle{}}l@{}}#1\end{array}$}}
\newcommand{\lycee@eqmesure}[1]{%
  \ifdim\ht#1>\lycee@eqh\setlength{\lycee@eqh}{\ht#1}\fi
  \ifdim\dp#1>\lycee@eqd\setlength{\lycee@eqd}{\dp#1}\fi}
\newcommand{\lycee@eqtrait}{\hspace{1.2em}\tikz[baseline=0pt]\draw[dash pattern=on 2pt off 2pt,
  line width=0.5pt] (0,\lycee@eqh) -- (0,-\lycee@eqd);\hspace{1.2em}}
\newcommand{\eqproduit}[2]{\par\smallskip\noindent C'est une équation produit :
  $A\times B=0 \iff A=0$ ou $B=0$.\par\nopagebreak\smallskip
  \begingroup\setlength{\lycee@eqh}{0pt}\setlength{\lycee@eqd}{0pt}%
  \setbox\lycee@eqa=\lycee@eqcol{#1}\setbox\lycee@eqb=\lycee@eqcol{#2}%
  \lycee@eqmesure\lycee@eqa\lycee@eqmesure\lycee@eqb%
  \noindent\hspace*{\retraitcalculs}\usebox\lycee@eqa\lycee@eqtrait\usebox\lycee@eqb\endgroup\par\smallskip}
\newcommand{\eqproduittrois}[3]{\par\smallskip\noindent C'est une équation produit :
  $A\times B\times C=0 \iff A=0$ ou $B=0$ ou $C=0$.\par\nopagebreak\smallskip
  \begingroup\setlength{\lycee@eqh}{0pt}\setlength{\lycee@eqd}{0pt}%
  \setbox\lycee@eqa=\lycee@eqcol{#1}\setbox\lycee@eqb=\lycee@eqcol{#2}\setbox\lycee@eqc=\lycee@eqcol{#3}%
  \lycee@eqmesure\lycee@eqa\lycee@eqmesure\lycee@eqb\lycee@eqmesure\lycee@eqc%
  \noindent\hspace*{\retraitcalculs}\usebox\lycee@eqa\lycee@eqtrait\usebox\lycee@eqb\lycee@eqtrait\usebox\lycee@eqc\endgroup
  \par\smallskip}

\tikzset{grille/.style={coulGrille,line width=0.4pt}}
\newcommand{\quadrillage}[4]{\draw[grille,step=1] (#1,#2) grid (#3,#4);}
\tikzset{
  croix/.style={cross out,draw,line width=0.6pt,inner sep=0pt,minimum size=4pt},
  croixdroite/.style={croix,rotate=45,line width=0.8pt},
}
\newcommand{\pt}[3]{\path (#1) node[croix]{} node[#2,font=\small,inner sep=2.5pt]{$#3$};}

\newcommand{\lycee@repere}[4]{%
  \draw[grille,step=1] (#1,#2) grid (#3,#4);
  \draw[-{Stealth[length=2mm]},thick] (#1,0)--({#3+0.4},0) node[below left=-1pt]{$x$};
  \draw[-{Stealth[length=2mm]},thick] (0,#2)--(0,{#4+0.4}) node[below left=-1pt]{$y$};
  \node[below left,font=\scriptsize,inner sep=1.5pt] at (0,0) {$0$};}
\newcommand{\lycee@gradx}[1]{\foreach \k/\l in {#1}{%
  \draw (\k,0.07)--(\k,-0.07) node[below,font=\scriptsize,inner sep=1.5pt]{$\l$};}}
\newcommand{\lycee@grady}[1]{\foreach \k/\l in {#1}{%
  \draw (0.07,\k)--(-0.07,\k) node[left,font=\scriptsize,inner sep=1.5pt]{$\l$};}}
\newcommand{\lycee@intff}[2]{\left[#1\,;#2\right]}
\newcommand{\lycee@intfo}[2]{\left[#1\,;#2\right[}
\newcommand{\lycee@intof}[2]{\left]#1\,;#2\right]}
\newcommand{\lycee@intoo}[2]{\left]#1\,;#2\right[}
\newcommand{\lycee@ens}[1]{\left\{#1\right\}}
\AtBeginDocument{%
  \providecommand{\repere}{\lycee@repere}%
  \providecommand{\gradx}{\lycee@gradx}%
  \providecommand{\grady}{\lycee@grady}%
  \providecommand{\Cf}{\mathcal{C}\sb{f}}%
  \providecommand{\Cg}{\mathcal{C}\sb{g}}%
  \providecommand{\intff}{\lycee@intff}%
  \providecommand{\intfo}{\lycee@intfo}%
  \providecommand{\intof}{\lycee@intof}%
  \providecommand{\intoo}{\lycee@intoo}%
  \providecommand{\ens}{\lycee@ens}}

\newlist{questions}{enumerate}{3}
\setlist[questions,1]{label=\arabic*\textdegree),leftmargin=*,itemsep=2pt,topsep=3pt}
\setlist[questions,2]{label=\alph*),leftmargin=*,itemsep=1pt,topsep=2pt}
\setlist[questions,3]{label=\roman*),leftmargin=*,itemsep=1pt,topsep=1pt}
\setlist[itemize]{label=\textbullet,leftmargin=*,itemsep=2pt,topsep=3pt}

\newcommand{\besoinplace}[1]{\par\penalty\@M\begingroup
  \dimen@=#1\relax\dimen@ii=\pagegoal\advance\dimen@ii-\pagetotal
  \ifdim\dimen@>\dimen@ii\ifdim\dimen@ii>\z@\vfil\fi\break\fi\endgroup}

\newcommand{\lycee@pied}{\thepage/\pageref{LastPage}}
\newcommand{\entete}[2]{%
  \pagestyle{fancy}\fancyhf{}%
  \fancyhead[L]{\footnotesize\color{coulPart}\textbf{#1}}%
  \fancyhead[R]{\footnotesize\color{coulPart}\textbf{#2}}%
  \fancyfoot[C]{\footnotesize\lycee@pied}%
  \renewcommand{\headrulewidth}{0.4pt}%
  \renewcommand{\headrule}{\hbox to\headwidth{%
    \color{coulPart!40}\leaders\hrule height \headrulewidth\hfill}}}

\RequirePackage{ccicons}
\newcommand{\lycee@licence}{{\scriptsize\color{gray}%
  \copyright\ Luc Destombe --- \ccbyncsa\ CC BY-NC-SA 4.0}}
\newif\iflycee@licence \lycee@licencetrue
\newcommand{\sanslicence}{\lycee@licencefalse}
\AtBeginDocument{\iflycee@licence\AddToHookNext{shipout/foreground}{%
  \put(0,0){\rlap{\kern\dimexpr1in+\hoffset+\oddsidemargin+\textwidth\relax
    \llap{\raisebox{-\dimexpr1in+\voffset+\topmargin+\headheight+\headsep
      +\textheight+\footskip\relax}{\lycee@licence}}}}}\fi}

\newcommand{\lycee@site}{faire-des-maths.fr}
\newcommand{\lycee@version}{{\scriptsize\color{gray}\lycee@site
  \ifdefined\lyceeversion\ --- Version~\lyceeversion\fi}}
\AtBeginDocument{\iflycee@licence\AddToHookNext{shipout/foreground}{%
  \put(0,0){\rlap{\kern\dimexpr1in+\hoffset+\oddsidemargin\relax
    \raisebox{-\dimexpr1in+\voffset+\topmargin+\headheight+\headsep
      +\textheight+\footskip\relax}{\lycee@version}}}}\fi}

\newcommand{\formulesserrees}{%
  \setlength{\abovedisplayskip}{3pt}\setlength{\belowdisplayskip}{3pt}%
  \setlength{\abovedisplayshortskip}{2pt}\setlength{\belowdisplayshortskip}{3pt}}

\setlength{\parindent}{0pt}

\RequirePackage[implicit=false,hidelinks,pdfencoding=auto,psdextra,bookmarksopen,
  bookmarksopenlevel=1,pdfcreator={},pdfproducer={}]{hyperref}
\RequirePackage{bookmark}
\hypersetup{pdfauthor={Luc Destombe},pdfkeywords={CC BY-NC-SA 4.0}}
\newcounter{lycee@signet}
\newcommand{\signet}[3][]{\stepcounter{lycee@signet}%
  \edef\lycee@dest{\if\relax\detokenize{#1}\relax
    signet.\arabic{lycee@signet}\else#1\fi}%
  \hypertarget{\lycee@dest}{}\bookmark[level=#2,dest=\lycee@dest]{#3}}
\pdfstringdefDisableCommands{%
  \def\R{R}\def\N{N}\def\Z{Z}\def\Q{Q}\def\vect#1{#1}\def\Cf{Cf}\def\Cg{Cg}\def\fois#1{×#1}%
  \def\barre#1{#1}\def\intervalle#1#2{[#1;#2]}\def\intff#1#2{[#1;#2]}%
  \def\textsuperscript#1{#1}\def\dfrac#1#2{#1/#2}\def\frac#1#2{#1/#2}%
  \def\vec#1{#1⃗}}%
\RequirePackage[official]{eurosym}

\geometry{top=0.8cm,bottom=1.3cm,left=1cm,right=1cm,footskip=0.6cm}

\pagestyle{fancy}
\fancyhf{}
\renewcommand{\headrulewidth}{0pt}
\fancyfoot[C]{\footnotesize Page \thepage{} / \pageref{LastPage}}

\newcommand{\titrefiche}[2]{%
  \begin{center}
    \Large\textbf{#1}\\[2pt]
    \large\textbf{#2}\\[2pt]
  \end{center}
  \vspace{0.10cm}}

\setlength{\columnseprule}{0.4pt}
\setlength{\columnsep}{15pt}
\renewcommand{\columnseprulecolor}{\color{black!45}}
\setlength{\multicolsep}{2pt}

\newcounter{partiefiche}
\renewcommand{\thepartiefiche}{\Roman{partiefiche}}
\newif\ifapresbandeau
\newcommand{\lycee@bandeau}[1]{%
  \par\penalty-600
  \vspace{0.08cm}%
  \noindent\colorbox{coulPart}{%
    \parbox{\dimexpr\linewidth-2\fboxsep\relax}{%
      \color{white}\bfseries\raggedright\strut#1\strut}}%
  \nopagebreak\par\nopagebreak
  \vspace{0.06cm}\nopagebreak
  \apresbandeautrue}
\newcommand{\partiefiche}[1]{%
  \refstepcounter{partiefiche}%
  \lycee@bandeau{\signet{1}{\thepartiefiche{} --- #1}\thepartiefiche{} --- #1}}
\newcommand{\partiefichesansnum}[1]{\lycee@bandeau{\signet{1}{#1}#1}}

\newcounter{exo}
\newlength{\espaceexo}\setlength{\espaceexo}{7pt}
\newlength{\espacetitre}\setlength{\espacetitre}{2pt}
\newcommand{\exo}[1][\relax]{%
  \par
  \ifapresbandeau\apresbandeaufalse\else\penalty-150\addvspace{\espaceexo}\fi
  \refstepcounter{exo}%
  \noindent\signet[exercice-\arabic{exo}]{2}{Exercice \arabic{exo}}%
  \textbf{\color{coulExo}\underline{Exercice \theexo}}%
  \ifx\relax#1\relax\else\textbf{ : #1}\fi
  \par\nopagebreak\vspace{\espacetitre}\nopagebreak
  \ignorespaces}

\setlist[enumerate]{nosep,leftmargin=*,itemsep=2pt,topsep=2pt}
\setlist[itemize]{nosep,leftmargin=*,topsep=2pt}
\newcommand{\choix}[3]{\par\hspace*{1em}%
  \makebox[0.3\linewidth][l]{a) #1}\makebox[0.3\linewidth][l]{b) #2}c) #3}
\newcommand{\deux}[2]{\par\hspace*{0.6em}%
  \makebox[0.48\linewidth][l]{#1}#2}
\newcommand{\trois}[3]{\par\hspace*{0.6em}%
  \makebox[0.32\linewidth][l]{#1}\makebox[0.32\linewidth][l]{#2}#3}

  \renewcommand{\exo}[1][\relax]{%
    \ifapresbandeau\apresbandeaufalse\else\par\penalty-150\fi
    \refstepcounter{exo}%
    \vspace{0.05cm}%
    \noindent\signet[exercice-\arabic{exo}]{2}{Exercice \arabic{exo}}%
    \textbf{\color{coulExo}\underline{Exercice \theexo}}%
    \ifx\relax#1\relax\else\textbf{ : #1}\fi%
    \par\nopagebreak
    \ignorespaces}
  \newcommand{\rappel}[1]{%
    \par\vspace{0.06cm}\nopagebreak
    \noindent\fcolorbox{black!35}{black!5}{%
      \parbox{\dimexpr\linewidth-2\fboxsep-2\fboxrule\relax}{%
        \small\textbf{Rappel.} #1}}%
    \par\vspace{0.06cm}\nopagebreak}
  \setlist[enumerate]{nosep,leftmargin=*}
  \setlist[itemize]{nosep,leftmargin=*}
  \setlength{\multicolsep}{3pt plus 1pt minus 1pt}
  \setlength{\premulticols}{10pt}
  \setlength{\postmulticols}{10pt}
  \newlist{expr}{enumerate}{1}
  \setlist[expr]{label={\Alph*\ $=$},nosep,leftmargin=*,itemsep=0pt}

\renewenvironment{center}{\par\addvspace{3pt}\centering}{\par\addvspace{3pt}}
\formulesserrees
\AtBeginDocument{\formulesserrees}
\clubpenalty=10000
\widowpenalty=10000
\displaywidowpenalty=10000
\brokenpenalty=10000
\setlength{\parskip}{0pt}

\RequirePackage{ragged2e}
\setlength{\RaggedRightRightskip}{0pt plus 1fil}
\AtBeginDocument{\RaggedRight\binoppenalty=10000 \relpenalty=10000 }
\g@addto@macro\@parboxrestore{\RaggedRight}
\makeatother
\geometry{top=0.8cm,bottom=1.3cm,left=1cm,right=1cm,footskip=0.7cm,headheight=12pt}

\raggedcolumns
\newcommand{\qcm}[4]{\par\makebox[0.24\linewidth][l]{a) #1}%
  \makebox[0.24\linewidth][l]{b) #2}\makebox[0.24\linewidth][l]{c) #3}d) #4}

\usepackage{tkz-tab}
\definecolor{coulCourbe}{HTML}{C0392B}
\definecolor{coulCourbeB}{HTML}{1F4E79}

\begin{document}

\titrefiche{1\textsuperscript{re} STMG --- Fonctions et pourcentages}%
  {Fiche d'exercices du chapitre}

\begin{multicols}{2}

\partiefiche{Tableau de valeurs, courbe et lecture graphique}

\exo[Première lecture graphique]
Voici la courbe d'une fonction $f$ définie sur $\intff{-2}{4}$.
\begin{center}
\begin{tikzpicture}[x=0.4cm,y=0.4cm]
  \repere{-3}{-5}{5}{5}
  \gradx{-2,-1,1,2,3,4} \grady{-4,-3,-2,-1,1,2,3,4}
  \draw[coulCourbe,line width=1.1pt,domain=-2:4,samples=60] plot (\x,{-\x*\x+2*\x+3});
  \node[coulCourbe,font=\footnotesize] at (3.8,2.6) {$\Cf$};
\end{tikzpicture}
\end{center}
Compléter par lecture graphique.
\begin{enumerate}[label=\alph*)]
  \item $f(0)=\dots$ \quad et \quad $f(2)=\dots$
  \item L'image de $-1$ par $f$ est $\dots$
  \item $4$ a pour antécédent $\dots$
  \item Les antécédents de $0$ par $f$ sont $\dots$ et $\dots$
  \item Le nombre $5$ a-t-il un antécédent par $f$ ? Justifier.
\end{enumerate}

\exo[Lire un tableau de valeurs]
Voici un tableau de valeurs d'une fonction $f$.
\begin{center}\small
\renewcommand{\arraystretch}{1.25}
\begin{tabular}{|c|c|c|c|c|c|}
\hline
$x$ & $-2$ & $-1$ & $0$ & $1$ & $2$\\
\hline
$f(x)$ & $4$ & $2$ & $-1$ & $-3$ & $1$\\
\hline
\end{tabular}
\end{center}
\begin{enumerate}[label=\arabic*)]
  \item Donner l'image de $-1$, de $1$, puis de $-2$.
  \item Le nombre $1$ est-il un antécédent de $-3$ ?
  \item Donner un antécédent de $1$.
\end{enumerate}

\exo[Lecture graphique complète]
Voici la courbe d'une fonction $f$.
\begin{center}
\begin{tikzpicture}[x=0.4cm,y=0.4cm]
  \repere{-5}{-3}{6}{3}
  \gradx{-4,-3,-2,-1,1,2,3,4,5} \grady{-2,-1,1,2}
  \draw[coulCourbe,line width=1.1pt]
    (-4,-2) .. controls (-3,-1) and (-2,0) .. (-1,1)
    .. controls (-0.667,1.333) and (-0.333,2) .. (0,2)
    .. controls (0.333,2) and (0.667,1.333) .. (1,1)
    .. controls (1.333,0.667) and (1.667,0.333) .. (2,0)
    .. controls (2.333,-0.333) and (2.667,-1) .. (3,-1)
    .. controls (3.333,-1) and (3.667,-0.333) .. (4,0)
    .. controls (4.333,0.333) and (4.667,0.667) .. (5,1);
  \path[coulCourbe] (-4,-2) node[croix]{} (5,1) node[croix]{};
  \node[coulCourbe,font=\footnotesize] at (5.4,2) {$\Cf$};
\end{tikzpicture}
\end{center}
Déterminer, par lecture graphique :
\begin{multicols}{2}
\begin{enumerate}[label=\alph*)]
  \item l'intervalle sur lequel $f$ est définie ;
  \item l'image de $3$ ;
  \item l'image de $-1$ ;
  \item $f(1)$ et $f(-4)$ ;
  \item les antécédents de $2$ ;
  \item les antécédents de $1$ ;
  \item les antécédents de $0$ ;
  \item les antécédents de $-1$ ;
  \item les antécédents de $3$ ;
  \item le maximum et le minimum de $f$, et où ils sont atteints.
\end{enumerate}
\end{multicols}

\partiefiche{Tableau de variations d'une fonction}

\exo[De la courbe au tableau]
La fonction $f$ est définie sur $\intff{-3}{5}$ par sa courbe ci-dessous. Dresser son
tableau de variations, puis donner son maximum et son minimum.
\begin{center}
\begin{tikzpicture}[x=0.4cm,y=0.4cm]
  \repere{-4}{-3}{6}{5}
  \gradx{-3,-2,-1,1,2,3,4,5} \grady{-2,-1,1,2,3,4}
  \draw[coulCourbe,line width=1.1pt]
    (-3,3) .. controls (-2,0.5) and (0,-2) .. (1,-2)
           .. controls (2,-2) and (4,1.5) .. (5,4);
  \foreach \p in {(-3,3),(1,-2),(5,4)} \path[coulCourbe] \p node[croix]{};
  \node[coulCourbe,font=\footnotesize] at (5.4,2.6) {$\Cf$};
\end{tikzpicture}
\end{center}

\exo[Décrire puis résumer]
Voici la courbe d'une fonction $f$ définie sur $\intff{-3}{4}$.
\begin{center}
\begin{tikzpicture}[x=0.4cm,y=0.4cm]
  \repere{-4}{-4}{5}{3}
  \gradx{-3,-2,-1,1,2,3,4} \grady{-3,-2,-1,1,2}
  \draw[coulCourbe,line width=1.1pt]
    (-3,-1) .. controls (-2.33,0.4) and (-1.67,2) .. (-1,2)
            .. controls (0,2) and (1,-3) .. (2,-3)
            .. controls (2.67,-3) and (3.33,-1.6) .. (4,0);
  \foreach \p in {(-3,-1),(-1,2),(2,-3),(4,0)} \path[coulCourbe] \p node[croix]{};
  \node[coulCourbe,font=\footnotesize] at (4.4,1.2) {$\Cf$};
\end{tikzpicture}
\end{center}
\begin{enumerate}[label=\arabic*)]
  \item Décrire par des phrases les variations de $f$.
  \item Dresser le tableau de variations de $f$.
\end{enumerate}

\exo[Lire un tableau de variations]
Voici le tableau de variations d'une fonction $g$.
\begin{center}
\begin{tikzpicture}[scale=0.8]
  \tkzTabInit[lgt=1.4,espcl=1.5]{$x$/1,$g$/2}{$-5$,$-2$,$1$,$4$}
  \tkzTabVar{+/$3$,-/$-2$,+/$4$,-/$0$}
\end{tikzpicture}
\end{center}
\begin{enumerate}[label=\arabic*)]
  \item Sur quel intervalle $g$ est-elle définie ?
  \item Décrire par des phrases les variations de $g$.
  \item Donner le maximum et le minimum de $g$, et où ils sont atteints.
  \item Tracer dans un repère une courbe qui peut représenter $g$.
\end{enumerate}

\exo[Résolution graphique]
Voici la courbe d'une fonction $f$ définie sur $\intff{0}{10}$.
\begin{center}
\begin{tikzpicture}[x=0.4cm,y=0.4cm]
  \repere{-1}{-3}{11}{4}
  \gradx{1,2,3,4,5,6,7,8,9,10} \grady{-2,-1,1,2,3}
  \draw[coulCourbe,line width=1.1pt]
    (0,2) .. controls (1,1) and (2,0) .. (3,-1)
    .. controls (3.333,-1.333) and (3.667,-2) .. (4,-2)
    .. controls (4.333,-2) and (4.667,-1.333) .. (5,-1)
    .. controls (6,0) and (7,1) .. (8,2)
    .. controls (8.333,2.333) and (8.667,3) .. (9,3)
    .. controls (9.333,3) and (9.667,2.333) .. (10,2);
  \path[coulCourbe] (0,2) node[croix]{} (10,2) node[croix]{};
  \node[coulCourbe,font=\footnotesize] at (6.4,1.9) {$\Cf$};
\end{tikzpicture}
\end{center}
\begin{enumerate}[label=\arabic*)]
  \item Résoudre graphiquement :
        \begin{multicols}{3}
        \begin{enumerate}[label=\alph*)]
          \item $f(x)=2$
          \item $f(x)=1$
          \item $f(x)=-2$
          \item $f(x)=4$
          \item $f(x)=0$
        \end{enumerate}
        \end{multicols}
  \item Combien l'équation $f(x)=-1$ a-t-elle de solutions ?
  \item Résoudre graphiquement, en donnant les solutions sous forme d'intervalles :
        \begin{multicols}{2}
        \begin{enumerate}[label=\alph*)]
          \item $f(x)>-2$
          \item $f(x)<1$
          \item $f(x)\ge 1$
          \item $f(x)\ge 0$
        \end{enumerate}
        \end{multicols}
\end{enumerate}

\partiefiche{Appliquer un pourcentage}

\exo[Calcul mental]
Calculer de tête.
\trois{a) $10\,\%$ de $420$}{b) $25\,\%$ de $60$}{c) $50\,\%$ de $48$}
\trois{d) $1\,\%$ de $900$}{e) $5\,\%$ de $80$}{f) $75\,\%$ de $20$}
\trois{g) $20\,\%$ de $35$}{h) $30\,\%$ de $250$}{i) $40\,\%$ de $15$}

\exo[Trois écritures d'un même nombre]
Recopier et compléter le tableau.
\begin{center}
\renewcommand{\arraystretch}{2.4}
\begin{tabular}{|l|c|c|c|c|c|}
\hline
Fraction & $\dfrac{3}{5}$ & & & $\dfrac{1}{4}$ & \\
\hline
Décimal & & $0{,}09$ & & & $1{,}5$\\
\hline
Pourcentage & & & $45\,\%$ & & \\
\hline
\end{tabular}
\end{center}

\exo[Retrouver le total]
Pour chaque situation, calculer l'effectif total.
\begin{enumerate}[label=\alph*)]
  \item $30$ chevaux représentent $5\,\%$ des animaux d'une ferme.
  \item $9$ élèves représentent $30\,\%$ d'une classe.
  \item $8$~€ représentent $10\,\%$ du prix d'un jeu.
  \item $60$ adhérents représentent $75\,\%$ d'un club.
\end{enumerate}

\exo[Augmenter ou diminuer]
Calculer chaque nouvelle valeur, en calculant d'abord le montant de la hausse ou de la
baisse.
\begin{enumerate}[label=\alph*)]
  \item Un article à $60$~€ augmente de $15\,\%$.
  \item Un ordinateur à $900$~€ est soldé à $-20\,\%$.
  \item Un salaire de $2\,000$~€ augmente de $3\,\%$.
  \item Une commune de $6\,000$ habitants perd $5\,\%$ de sa population.
\end{enumerate}

\exo[Prix hors taxe, prix TTC]
Une trottinette coûte $350$~€ hors taxe (HT). On lui ajoute la TVA, de $20\,\%$.
\begin{enumerate}[label=\alph*)]
  \item Calculer le montant de la TVA, puis le prix toutes taxes comprises (TTC).
  \item En soldes, le prix TTC baisse de $10\,\%$. Calculer le prix soldé.
\end{enumerate}

\partiefiche{Calculer un pourcentage}

\exo[Proportions]
Écrire chaque proportion en pourcentage.
\deux{a) $12$ élèves sur $40$}{b) $9$ admis sur $20$ candidats}
\deux{c) $33$ garçons sur $150$ élèves}{d) $80$ clients sur $500$}

\exo[Taux d'évolution]
Calculer le taux d'évolution de chaque grandeur, en pourcentage.
\begin{enumerate}[label=\alph*)]
  \item Un prix passe de $150$~€ à $180$~€.
  \item Un stock passe de $400$~kg à $340$~kg.
  \item Un loyer passe de $800$~€ à $720$~€.
  \item Un club passe de $60$ à $75$ adhérents.
\end{enumerate}

\exo[Plus ou moins de\ldots{} ?]
Répondre sans calculer le taux d'évolution exact.
\begin{enumerate}[label=\alph*)]
  \item Un club passe de $700$ à $620$ adhérents. A-t-il perdu plus de $10\,\%$ de ses
        adhérents ?
  \item Une forêt passe de $1\,500$ à $1\,620$ arbres. A-t-elle gagné plus de
        $10\,\%$ d'arbres ?
  \item Une ville passe de $3\,000$ à $2\,880$ habitants. A-t-elle perdu plus de
        $5\,\%$ de ses habitants ?
\end{enumerate}

\exo[Au tableur]
Une feuille de calcul donne le nombre d'entrées d'une piscine.
\begin{center}
\renewcommand{\arraystretch}{1.25}
\small
\begin{tabular}{|c|l|c|c|c|c|}
\hline
 & \multicolumn{1}{c|}{A} & B & C & D & E\\
\hline
1 & Année & 2022 & 2023 & 2024 & 2025\\
\hline
2 & Entrées & $5\,000$ & $5\,500$ & $4\,950$ & $5\,445$\\
\hline
3 & Taux d'évolution & \large$\times$ & & & \\
\hline
\end{tabular}
\end{center}
\rappel{Dans une formule, \texttt{B2} désigne le nombre de la cellule B2. En recopiant
une formule vers la droite, \texttt{B2} devient \texttt{C2}, puis \texttt{D2}\ldots}
\begin{enumerate}[label=\alph*)]
  \item Calculer le taux d'évolution du nombre d'entrées entre 2022 et 2023.
  \item Parmi ces formules, laquelle saisir en C3 pour obtenir, par recopie vers la
        droite, les taux d'évolution d'une année sur l'autre ?
        \par\hspace*{0.6em}\makebox[0.45\linewidth][l]{\texttt{=(C2-B2)/C2}}\texttt{=C2-B2/B2}
        \par\hspace*{0.6em}\makebox[0.45\linewidth][l]{\texttt{=(C2-B2)/B2}}\texttt{=(B2-C2)/B2}
  \item Quelle formule contient alors la cellule D3 ?
\end{enumerate}

\partiefiche{Coefficient multiplicateur}

\exo[Du taux au coefficient]
Donner le coefficient multiplicateur de chaque évolution.
\trois{a) $+6\,\%$}{b) $-6\,\%$}{c) $+30\,\%$}
\trois{d) $-30\,\%$}{e) $+65\,\%$}{f) $-12\,\%$}
\trois{g) $+4\,\%$}{h) $-45\,\%$}{i) $+200\,\%$}

\exo[Du coefficient au taux]
Donner l'évolution, en pourcentage, associée à chaque coefficient.
\trois{a) $\times 1{,}09$}{b) $\times 0{,}91$}{c) $\times\frac{4}{5}$}
\trois{d) $\times 0{,}7$}{e) $\times 3$}{f) $\times 0{,}98$}

\exo[Appliquer une évolution]
Calculer chaque nouvelle valeur en une seule multiplication.
\begin{enumerate}[label=\alph*)]
  \item $300$~€ augmentés de $15\,\%$.
  \item $600$~€ diminués de $35\,\%$.
  \item $50$~€ augmentés de $120\,\%$.
  \item $2\,000$ habitants diminués de $1{,}5\,\%$.
\end{enumerate}

\exo[Le bon calcul]
Pour chaque question, une seule réponse est exacte.
\begin{enumerate}[label=\arabic*.]
  \item Un article à $500$~€ baisse de $12\,\%$. Son nouveau prix est :
        \qcm{$500-0{,}12$}{$500\times 0{,}88$}{$500\times 1{,}12$}{$500\times 0{,}12$}
  \item Un article à $150$~€ augmente de $30\,\%$. Son nouveau prix est :
        \qcm{$150+0{,}3$}{$150\times\dfrac{30}{100}$}{$150\times 1{,}3$}{$150\times 0{,}7$}
\end{enumerate}

\partiefiche{Évolutions successives, évolution réciproque}

\exo[Taux global]
Calculer le taux d'évolution global.
\begin{enumerate}[label=\alph*)]
  \item Une hausse de $50\,\%$, puis une hausse de $20\,\%$.
  \item Une baisse de $30\,\%$, puis une hausse de $30\,\%$.
  \item Une hausse de $100\,\%$, puis une baisse de $50\,\%$.
  \item Deux baisses successives de $20\,\%$.
  \item Une hausse de $60\,\%$, puis une baisse de $25\,\%$.
\end{enumerate}

\exo[Le prix de l'électricité]
Un journaliste annonce : « Une hausse de $10\,\%$ en février, puis une autre de
$10\,\%$ en août, augmenteront le prix de l'électricité de $20\,\%$ sur l'année. »
A-t-il raison ?

\exo[Une même évolution répétée]
Donner chaque fois le taux d'évolution global.
\begin{enumerate}[label=\alph*)]
  \item Une entreprise augmente sa production de $3\,\%$ par an. Par quel nombre la
        production est-elle multipliée en $3$ ans ? On donne
        $1{,}03^{3}\approx 1{,}09$.
  \item Les prix augmentent de $1\,\%$ par mois pendant un an. On donne
        $1{,}01^{12}\approx 1{,}13$. Quel est le taux d'évolution annuel ?
\end{enumerate}

\exo[Se compensent-elles ?]
Pour chaque cas, dire si la seconde évolution ramène à la valeur de départ, en
calculant le produit des coefficients multiplicateurs.
\deux{a) $-20\,\%$, puis $+25\,\%$}{b) $+50\,\%$, puis $-50\,\%$}
\deux{c) $-50\,\%$, puis $+100\,\%$}{d) $+300\,\%$, puis $-75\,\%$}
\deux{e) $-20\,\%$, puis $+20\,\%$}{f) $+10\,\%$, puis $-10\,\%$}

\exo[Quelle évolution compense ?]
Pour chaque question, une seule réponse est exacte : tester chaque réponse.
\begin{enumerate}[label=\arabic*.]
  \item Un prix est multiplié par $4$. Pour revenir au prix de départ, il faut une
        baisse de :
        \qcm{$400\,\%$}{$75\,\%$}{$25\,\%$}{$300\,\%$}
  \item Un prix augmente de $25\,\%$. Pour revenir au prix de départ, il faut une
        baisse de :
        \qcm{$25\,\%$}{$75\,\%$}{$20\,\%$}{$80\,\%$}
  \item Une population baisse de $50\,\%$. Pour revenir à son niveau, il faut une
        hausse de :
        \qcm{$100\,\%$}{$50\,\%$}{$200\,\%$}{$150\,\%$}
\end{enumerate}

\partiefiche{Fonctions affines : image, antécédent}

\exo[Affine ou non ?]
Dire si chaque fonction est affine ; si oui, donner $a$ et $b$.
\trois{a) $f(x)=5x+1$}{b) $g(x)=x^{2}-3$}{c) $h(x)=2-3x$}
\trois{d) $k(x)=\dfrac{2}{x}$}{e) $m(x)=1{,}5x$}{f) $n(x)=-4$}

\exo[Images]
Soit $f(x)=3x-4$ et $g(x)=-4x+6$.
\begin{enumerate}[label=\alph*)]
  \item Calculer $f(0)$, $f(2)$, $f(-1)$ et $f(0{,}5)$.
  \item Calculer $g(0)$, $g(2)$, $g(-3)$ et $g(1{,}5)$.
\end{enumerate}

\exo[Antécédents]
Calculer chaque antécédent par une équation.
\begin{enumerate}[label=\alph*)]
  \item $f(x)=3x+9$ : calculer l'antécédent de $15$, de $0$, puis de $-6$.
  \item $g(x)=-2x+8$ : calculer l'antécédent de $0$, puis de $14$.
  \item $h(x)=0{,}5x-2$ : résoudre $h(x)=3$.
\end{enumerate}

\exo[Deux tarifs]
Répondre à chaque question par un calcul.
\begin{enumerate}[label=\alph*)]
  \item Un dépanneur facture $40$~€ de déplacement et $1{,}50$~€ par km. Exprimer le
        prix $T(x)$ pour $x$~km. Quelle distance a-t-il parcourue pour $70$~€ ?
  \item Une salle de réception se loue $200$~€, plus $15$~€ par invité. Exprimer le
        prix $S(x)$ pour $x$ invités. Calculer $S(40)$. Combien d'invités pour
        $650$~€ ?
\end{enumerate}

\partiefiche{Trouver une fonction affine}

\exo[Mettre en équation]
Exprimer chaque grandeur par une fonction affine de $x$.
\begin{enumerate}[label=\alph*)]
  \item Une location de kayak coûte $8$~€, plus $6$~€ par heure. $L(x)$ est le prix
        pour $x$ heures.
  \item Un réservoir de $60$~L consomme $5$~L tous les $100$~km. $R(x)$ est la quantité
        restante après $x$~km.
  \item Une vendeuse touche $1\,400$~€ par mois, plus $4\,\%$ de ses ventes. $S(x)$ est
        son salaire pour $x$~€ de ventes.
\end{enumerate}

\exo[Coefficient directeur]
Calculer le coefficient directeur de la droite $(AB)$.
\deux{a) $A(1\,;2)$ et $B(4\,;11)$}{b) $A(0\,;6)$ et $B(2\,;2)$}
\deux{c) $A(-2\,;3)$ et $B(-4\,;7)$}{d) $A(-3\,;-2)$ et $B(1\,;6)$}

\exo[Déterminer une fonction affine]
Pour chaque cas, déterminer la fonction affine $f$ dont la droite passe par $A$ et $B$.
\deux{a) $A(1\,;5)$ et $B(3\,;11)$}{b) $A(0\,;-4)$ et $B(2\,;2)$}
\deux{c) $A(-2\,;3)$ et $B(2\,;3)$}{d) $A(-4\,;-5)$ et $B(2\,;4)$}

\partiefiche{Tracer une fonction affine}

\exo[Sens de variation]
Donner le sens de variation de chaque fonction affine.
\trois{a) $f(x)=4-x$}{b) $g(x)=3x+1$}{c) $k(x)=-2$}
\trois{d) $p(x)=-5x$}{e) $h(x)=-0{,}2x+3$}{}
\par\hspace*{0.6em}f)~$m(x)=0{,}3x-1$

\exo[Tracer des droites]
Dans un même repère orthonormé, tracer les droites qui représentent :
\trois{$f(x)=x+1$}{$g(x)=-3x+2$}{$h(x)=0{,}5x-2$}

\exo[Location de paddles]
Un paddle se loue $4$~€, plus $4$~€ par heure. On note $L(x)$ le prix pour $x$
heures, avec $x$ entre $0$ et $8$.
\begin{enumerate}[label=\alph*)]
  \item Exprimer $L(x)$ en fonction de $x$.
  \item Calculer $L(0)$ et $L(8)$.
  \item Tracer la droite qui représente $L$ (en abscisse, $1$~cm pour $1$ heure ; en
        ordonnée, $1$~cm pour $5$~€).
  \item Lire graphiquement la durée de location pour $24$~€, puis vérifier par le
        calcul.
\end{enumerate}

\exo[Proportionnalité et fonctions linéaires]
Répondre à chaque question par un calcul.
\begin{enumerate}[label=\alph*)]
  \item Une boisson contient $10$~g de sucre pour $100$~mL. Exprimer la masse de sucre
        $f(x)$ dans $x$~mL de boisson, puis calculer $f(330)$.
  \item Un pouce vaut environ $2{,}5$~cm. Convertir $40$ pouces en centimètres, puis
        $75$~cm en pouces.
  \item Parmi $f(x)=-2x$ ; $g(x)=2x+3$ ; $h(x)=\dfrac{x}{4}$ et $k(x)=-1$, lesquelles
        sont linéaires ?
  \item La droite d'une fonction linéaire $\ell$ passe par le point $(5\,;15)$.
        Déterminer $\ell(x)$.
\end{enumerate}

\partiefiche{Lire l'équation d'une droite}

\exo[Lire une équation réduite]
\begin{minipage}[t]{0.36\linewidth}\raggedright\vspace{0pt}
Donner l'équation réduite de chaque droite, puis le sens de variation de la fonction
qu'elle représente.
\end{minipage}\hfill\begin{minipage}[t]{0.62\linewidth}\centering\vspace{0pt}
\begin{tikzpicture}[x=0.4cm,y=0.4cm]
  \quadrillage{-4}{-4}{5}{5}
  \draw[->,>=Stealth,black!70] (-4,0) -- (5.5,0) node[below,font=\tiny] {$x$};
  \draw[->,>=Stealth,black!70] (0,-4) -- (0,5.5) node[left,font=\tiny] {$y$};
  \foreach \i in {-4,...,-1,1,2,...,4} {\node[below,font=\tiny] at (\i,0) {$\i$};}
  \foreach \i in {-4,...,-1,1,2,...,4} {\node[left,font=\tiny] at (0,\i) {$\i$};}
  \node[below left,font=\tiny] at (0,0) {$0$};
  \begin{scope}
    \clip (-4,-4) rectangle (5,5);
    \draw[coulCourbe,thick] (-1,-4) -- (2,5);
    \draw[coulCourbeB,thick] (-3,5) -- (5,-3);
    \draw[coulCourbeC,thick] (-4,-0.333) -- (5,2.667);
    \draw[black!70,thick] (-4,-2) -- (5,-2);
  \end{scope}
  \node[coulCourbe,right,font=\small] at (2,4.4) {$(d_{1})$};
  \node[coulCourbeB,left,font=\small] at (-2.6,4.4) {$(d_{2})$};
  \node[coulCourbeC,above,font=\small] at (4.2,2.4) {$(d_{3})$};
  \node[black!70,above,font=\small] at (4.2,-2) {$(d_{4})$};
\end{tikzpicture}
\end{minipage}

\exo[Points d'une droite]
Soit $(d)$ la droite d'équation réduite $y=-2x+5$.
\begin{enumerate}[label=\alph*)]
  \item Les points $A(1\,;3)$, $B(3\,;1)$ et $C(-2\,;9)$ sont-ils sur $(d)$ ?
  \item Calculer l'ordonnée du point de $(d)$ d'abscisse $4$.
  \item Calculer l'abscisse du point de $(d)$ d'ordonnée $-7$.
\end{enumerate}

\exo[Équation réduite sans graphique]
Pour chaque question, une seule réponse est exacte.
\begin{enumerate}[label=\arabic*.]
  \item Une droite coupe l'axe des ordonnées en $2$ ; pour un déplacement horizontal
        de $3$, son déplacement vertical est $+1$. Son équation réduite est :
        \deux{a) $y=3x+2$}{b) $y=\dfrac{1}{3}x+2$}
        \deux{c) $y=x+\dfrac{1}{3}$}{d) $y=\dfrac{1}{3}x-2$}
  \item La droite qui passe par $(0\,;-3)$ et $(2\,;1)$ a pour équation réduite :
        \deux{a) $y=-3x+2$}{b) $y=0{,}5x-3$}
        \deux{c) $y=2x-3$}{d) $y=-2x-3$}
  \item La droite d'équation $y=4x-1$ :
        \par a) monte \quad b) descend \quad c) est horizontale
\end{enumerate}

\partiefiche{Comparer deux fonctions affines}

\exo[Équations]
Résoudre chaque équation.
\deux{a) $6x+5=4x+11$}{b) $3x-2=x+8$}
\deux{c) $2x+3=5x-9$}{d) $4x+1=-2x+4$}

\exo[Deux offres d'emploi]
Une commerciale hésite entre deux sociétés. Société A : $1\,200$~€ par mois, plus
$10\,\%$ des ventes. Société B : $800$~€ par mois, plus $20\,\%$ des ventes. On note
$x$ le montant des ventes du mois, en euros.
\begin{enumerate}[label=\alph*)]
  \item Exprimer les salaires $f(x)$ et $g(x)$ proposés par A et par B.
  \item Pour quel montant de ventes les deux salaires sont-ils égaux ?
  \item Elle pense vendre pour $5\,000$~€ par mois. Quelle société choisir ?
\end{enumerate}

\exo[Deux formules de piscine]
Formule A : $5$~€ l'entrée. Formule B : abonnement de $24$~€, puis $2$~€ l'entrée. Le
graphique donne le prix payé selon le nombre $x$ d'entrées.
\begin{center}
\begin{tikzpicture}[x=0.45cm,y=0.05cm]
  \draw[grille,xstep=1,ystep=10] (0,0) grid (12,60);
  \draw[->,>=Stealth,black!70] (0,0) -- (12.7,0) node[below,font=\tiny] {$x$};
  \draw[->,>=Stealth,black!70] (0,0) -- (0,66) node[left,font=\tiny] {€};
  \foreach \i in {1,...,12} {\node[below,font=\tiny] at (\i,0) {$\i$};}
  \foreach \v in {10,20,...,60} {\node[left,font=\tiny] at (0,\v) {$\v$};}
  \node[below left,font=\tiny] at (0,0) {$0$};
  \draw[coulCourbe,thick] (0,0) -- (12,60);
  \draw[coulCourbeB,thick] (0,24) -- (12,48);
  \node[coulCourbe,left,font=\small] at (10.6,54) {A};
  \node[coulCourbeB,below,font=\small] at (11.4,45) {B};
\end{tikzpicture}
\end{center}
\begin{enumerate}[label=\alph*)]
  \item Exprimer les prix $f(x)$ et $g(x)$ des formules A et B.
  \item Lire les coordonnées du point d'intersection des deux droites.
  \item Retrouver ce résultat par le calcul.
  \item Quelle formule choisir selon le nombre d'entrées ?
\end{enumerate}

\partiefiche{Fonction carré}

\exo[Images et antécédents par la fonction carré]
\begin{enumerate}[label=\arabic*)]
  \item Calculer le carré de $-4$, de $7$, de $-0{,}2$, puis de $20$.
  \item Donner les antécédents par la fonction carré de $25$, de $0$, de $-9$, puis
        de $11$.
\end{enumerate}

\exo[Lecture graphique avec la fonction carré]
On considère la fonction carré $f(x)=x^{2}$, pour $x$ dans $\intff{-2}{2}$.
\begin{enumerate}[label=\arabic*)]
  \item Compléter le tableau de valeurs.
        \begin{center}\small
        \setlength{\tabcolsep}{3.5pt}\renewcommand{\arraystretch}{1.4}
        \begin{tabular}{|c|c|c|c|c|c|c|c|c|c|}
        \hline
        $x$ & $-2$ & $-1{,}5$ & $-1$ & $-0{,}5$ & $0$ & $0{,}5$ & $1$ & $1{,}5$ & $2$\\
        \hline
        $f(x)$ & & & & & & & & & \\
        \hline
        \end{tabular}
        \end{center}
  \item Tracer la courbe de $f$ dans un repère orthonormé (unité : 1~cm).
  \item Résoudre graphiquement :\\
        a)~$x^{2}=1$ \quad b)~$x^{2}=4$ \quad c)~$x^{2}=-1$
  \item Résoudre graphiquement, en donnant les solutions sous forme d'intervalles :\\
        a)~$x^{2}\le 1$ \quad b)~$x^{2}>1$ \quad c)~$x^{2}<4$
\end{enumerate}

\exo[Équations avec des carrés]
Résoudre dans $\R$ :
\begin{multicols}{2}
\begin{enumerate}[label=\alph*)]
  \item $3x^{2}-8=x^{2}+10$
  \item $6x^{2}-1=8x^{2}-11$
  \item $2x^{2}+7=3$
\end{enumerate}
\end{multicols}

\partiefichesansnum{Exercices type épreuve anticipée}

\exo[Automatismes (QCM)]
Pour chaque question, une seule réponse est exacte.
\begin{enumerate}[label=\arabic*.]
  \item $30\,\%$ de $70$ est égal à :
        \qcm{$21$}{$210$}{$40$}{$2{,}1$}
  \item Un article à $100$~€ augmente de $20\,\%$, puis baisse de $20\,\%$. Il coûte :
        \qcm{$100$~€}{$96$~€}{$104$~€}{$80$~€}
  \item Lequel de ces nombres est inférieur à $0{,}2$ ?
        \qcm{$\dfrac{1}{4}$}{$25\,\%$}{$\dfrac{1}{8}$}{$0{,}3$}
  \item Une baisse de $20\,\%$ est compensée par une hausse de :
        \qcm{$20\,\%$}{$25\,\%$}{$80\,\%$}{$120\,\%$}
  \item La droite qui passe par $A(2\,;1)$ et $B(4\,;7)$ a pour coefficient
        directeur :
        \qcm{$3$}{$\dfrac{1}{3}$}{$2$}{$-3$}
  \item La fonction $f(x)=3x-4$ est :
        \par a) croissante \quad b) décroissante \quad c) constante
  \item L'antécédent de $13$ par $f(x)=2x+3$ est :
        \qcm{$29$}{$5$}{$8$}{$6{,}5$}
  \item Le point d'abscisse $3$ de la droite d'équation $y=-2x+4$ a pour ordonnée :
        \qcm{$-2$}{$10$}{$2$}{$-10$}
  \item Une distance de $x$ miles vaut $1{,}6x$ kilomètres. Pour $x=25$, cela fait :
        \qcm{$40$~km}{$26{,}6$~km}{$15{,}6$~km}{$4$~km}
\end{enumerate}

\exo[Les inscrits d'un club de judo]
Une feuille de calcul donne le nombre d'inscrits d'un club de judo.
\begin{center}
\renewcommand{\arraystretch}{1.25}
\small
\begin{tabular}{|c|l|c|c|c|c|}
\hline
 & \multicolumn{1}{c|}{A} & B & C & D & E\\
\hline
1 & Année & 2020 & 2021 & 2022 & 2023\\
\hline
2 & Rang $x$ & $0$ & $1$ & $2$ & $3$\\
\hline
3 & Inscrits & $180$ & $171$ & $166$ & $156$\\
\hline
4 & Taux d'évolution & \large$\times$ & & & \\
\hline
\end{tabular}
\end{center}
\begin{enumerate}[label=\arabic*.]
  \item Calculer le taux d'évolution du nombre d'inscrits entre 2020 et 2021.
  \item Quelle formule saisir en C4 pour obtenir, par recopie vers la droite, les
        taux d'évolution d'une année sur l'autre ?
  \item On modélise le nombre d'inscrits par $f(x)=-8x+180$, où $x$ est le rang de
        l'année. Donner le sens de variation de $f$ et interpréter $-8$.
  \item Selon ce modèle, combien d'inscrits y aura-t-il en 2028 ?
  \item À partir de quelle année y aura-t-il moins de $100$ inscrits ?
\end{enumerate}

\exo[Les soldes]
Un magasin fait une remise de $25\,\%$ sur un casque à $80$~€ et de $10\,\%$ sur une
enceinte à $120$~€.
\begin{enumerate}[label=\arabic*.]
  \item Calculer le prix soldé de chaque article.
  \item On achète les deux articles. Calculer le prix total payé, puis le pourcentage
        global de remise.
  \item Après les soldes, le prix de l'enceinte augmente de $10\,\%$. Revient-il à
        $120$~€ ?
\end{enumerate}

\end{multicols}
\end{document}
